\section{Conclusion and Future Work}
\label{sec:concl}

This paper presented AER-MQoS, a context-aware multi-objective RPL extension that unifies traffic-class-aware QoS differentiation, energy awareness, and link-reliability-informed routing within a single coherent framework. The protocol introduces a Multi-Criteria Score (MCS) driven by a dynamic context coefficient $\gamma$ that balances QoS responsiveness ($\alpha$) and energy preservation ($\beta$) according to application domain, traffic class, and residual energy. Application-layer weighted round-robin scheduling ($6{:}3{:}2{:}1$) mitigates inter-class interference above UDP without modifying RPL's control plane. The MCS is coupled to MRHOF-style hysteresis through an experimental objective function (OCP~8), and a bounded Q-learning specification for online $\gamma$ adjustment ($\pm15$ hundredths) is documented in Section~\ref{sec:ql} and left for future ablation.

The open-source implementation, simulation scripts, and raw CSV datasets are publicly available. Evaluation across four firmware variants on a 25-node Cooja topology ($N{=}4$ seeds, 1800~s) demonstrates correct four-way integration: all variants achieve $>93\%$ PDR under lossy UDGM ($\mathrm{success\_ratio}{=}0.85$) with per-seed ranges of $93.4$--$99.8\%$. Per-class PDR shows preliminary signs of QoS-aware rank shaping on some seeds (AER-MQoS $95.96\%$~C3 vs $79.52\%$~C0 on seed~20260603) alongside episodic sensitivity on the same seed, confirming the need for larger-scale campaigns. The primary contribution is an open, inspectable framework that makes the QoS--energy--trust trade-off explicit at both rank-computation and packet-scheduling levels.

\subsection*{Limitations}
This evaluation is simulation-only (Cooja/UDGM). NRE is a logical state variable, not Energest-calibrated. No adversarial attack scripts, ablations (A1--A3), or inferential tests appear in the committed dataset; per-lever contribution requires the planned campaigns of Section~\ref{sec:validation-plan}. Custom DIO/DAO TLVs for $\gamma$, NRE, and trust advertisement are not present in the evaluated binary (Section~\ref{sec:dio}). The overlapping PDR ranges confirm implementation correctness at the current traffic load; higher stress scenarios are needed to expose the protocol's differentiation envelope.

\subsection*{Future work}
The open firmware roadmap is organized in three phases. \textbf{Phase~I (short-term, simulation-only):} scale the Cooja matrix to $N{\geq}20$ seeds with Mann--Whitney confirmatory tests; execute ablations A1--A3 (WRR-off, learning-off, trust-neutral MCS); and stress-test with higher traffic loads ($2\times$~$C3$ rate), heterogeneous energy profiles, and lower success ratios ($\mathrm{success\_ratio}{\in}\{0.65,0.75\}$). \textbf{Phase~II (instrumentation):} Energest duty-cycle logging on Cooja to bound the gap between the NRE proxy and millijoule-scale consumption; characterize ROM/RAM footprints for MSP430-class targets (Sky, Zolertia RE-Mote). \textbf{Phase~III (extension):} scripted Cooja attack injections (selective forwarding, rank inflation) with the trust plane active; explore standard-compliant DIO/DAO TLV extensions for $\gamma$, NRE, and trust advertisement; and conduct parameter-frozen versus learning-enabled firmware comparisons.

These phases are designed to transition AER-MQoS from an integration-focused descriptive baseline to a confirmatory, multi-evidence protocol evaluation suitable for deployment-oriented venues.